% FCSG-Net B.Tech research project report.
% Built on the NIT Rourkela DissertationGuidelinesV4_5 template (copy kept in
% college_reference/). Needs the college bundle files listed in thesis/README.md
% (NITR.cls, FrontPages.tex, Ref/*.bst) before it will compile.
\documentclass[12pt,a4paper,oneside]{NITR}

\usepackage[hyphens]{url}
\usepackage[none]{hyphenat}
\usepackage{amsmath,amsfonts,mathdots,amssymb}  % template also loads yfonts (blackletter), unused here
\usepackage{graphicx}
\usepackage{xcolor}
\usepackage{tikz}
\usetikzlibrary{trees,shapes}
\usepackage{booktabs}
\usepackage{multirow}
\usepackage{longtable}
\usepackage{subcaption}
\usepackage{siunitx}            % aligned decimals in result tables
\sisetup{group-separator={,},group-minimum-digits=4}

% Report figures are already generated as vector PDFs by
% evaluation/build_technical_report.py. Reference them in place, do not copy.
\graphicspath{{figures/}}
% ===== begin macros.tex =====
% Shared macros. Keep this file small.

% Writer notes. Every \todo prints in red so nothing ships half-written.
% Before final submission: grep -rn '\\todo' thesis/ must return nothing.
\newcommand{\todo}[1]{\textcolor{red}{\textbf{[TODO:} #1\textbf{]}}}
% Where a number came from. Renders nothing, exists for reviewers of the .tex.
\newcommand{\evidence}[1]{}

\newcommand{\fcsg}{FCSG-Net}
\newcommand{\dB}{\,dB}
\newcommand{\expert}[1]{$E_{#1\times#1}$}   % \expert{7} -> E_{7x7}
% ===== end macros.tex =====


\onehalfspacing
\begin{document}
\sloppy
% ---------------------------------------------------------------- front matter
% ===== begin FrontPages.tex =====
% ======================================================================
% Front-page fields. STAND-IN layout: when the college FrontPages.tex
% arrives, copy these values into its fields and replace this file.
% ======================================================================
\reportTitle{FCSG-Net: Frequency-Band Spatial Expert Routing\\for Lightweight Composite Image Restoration}
\docType{Project Report}
\degree{Bachelor of Technology}
\dept{Computer Science and Engineering}
% Ascending order of roll number.
\authors{Akshat Mittal (123CS0954)\\Risav Rathi (123CS0957)\\Ramnath Thakur (123CS0965)}
\authorsInline{Akshat Mittal (123CS0954), Risav Rathi (123CS0957) and Ramnath Thakur (123CS0965)}
\newcommand{\supname}{Dr.\ Shyampada Mukherjee}
\supervisorOne{\supname}
\submitDate{November 2026}

\ackText{%
We thank our supervisor, \supname, for guiding us through this project.

We thank the Department of Computer Science and Engineering, National
Institute of Technology Rourkela, for the chance to do this work. Kaggle
gave us free GPU time. Without it we could not have trained four models for
200,000 updates each.

We thank the DIV2K authors for making their dataset public. We also thank
Kai Zhang and colleagues for releasing the FFDNet code and pretrained
weights, which let us check our implementation against a published result.

Last, we thank our families and friends for their patience and support.
}

\abstractText{%
Old photographs usually carry several kinds of damage at once. Optical blur,
lost resolution, sensor or film noise and lossy JPEG compression each hit a
different part of the image spectrum. We ask whether a very small
restoration network can make use of that. Our network, FCSG-Net, splits the
degraded image into low, mid and high frequency bands with fixed, soft radial
masks in the Fourier domain. Three shared convolutional experts, with
$7\times7$, $5\times5$ and $3\times3$ kernels, process every band. A small
gate picks the two best experts for each band at every pixel. The model has
97,088 parameters and costs 9.39 GFLOPs for a $256\times256$ input. We first
verified the implementation and reproduced the published FFDNet result to
within 0.01\,dB. We then trained FCSG-Net, DnCNN and two FFDNet variants from
scratch on the same synthetic composite degradation of DIV2K images, with one
schedule and one evaluation protocol. On 100 validation images FCSG-Net
reaches 27.24\,dB PSNR. DnCNN reaches 27.83\,dB and blind FFDNet 28.19\,dB,
and FCSG-Net uses 82.6\% fewer parameters than DnCNN. The low FLOP count did
not make it fast. On a Tesla T4 GPU it is the slowest of the four. The
learned routing varies across the image and follows the noise level, with
within-image correlations up to 0.81. Overriding the gate's choices at
inference costs 0.37\,dB on average. Frequency-band routing gives a small
model behaviour we can inspect. Within this budget, it does not make the
model more accurate or faster than simple convolutional baselines.
}

\keywords{image restoration, frequency decomposition, mixture of experts,
spatial routing, lightweight networks}
% ===== end FrontPages.tex =====
            % college file: title, names, roll nos, guide(s)

\coverPage
\titlePage
\certificateOfExamination

\certificateOneSupervisor
%\certificateTwoSupervisors

%\dedication                    % add \dedicationText in FrontPages.tex to use
\declaration
\acknowledgment
\abstract

\tableofcontents
\cleardoublepage
\addcontentsline{toc}{chapter}{List of Figures}
\listoffigures
\cleardoublepage
\addcontentsline{toc}{chapter}{List of Tables}
\listoftables
\cleardoublepage
% ===== begin chapters/abbreviations.tex =====
\chapter*{List of Abbreviations}
\addcontentsline{toc}{chapter}{List of Abbreviations}
\begin{tabular}{@{}ll@{}}
AMP    & Automatic Mixed Precision \\
AWGN   & Additive White Gaussian Noise \\
CNN    & Convolutional Neural Network \\
DIV2K  & DIVerse 2K resolution image dataset \\
FFT    & Fast Fourier Transform \\
FLOP   & Floating-Point Operation \\
GFLOP  & $10^9$ Floating-Point Operations \\
GPU    & Graphics Processing Unit \\
JPEG   & Joint Photographic Experts Group (compression) \\
LPIPS  & Learned Perceptual Image Patch Similarity \\
MAC    & Multiply-Accumulate Operation \\
MoE    & Mixture of Experts \\
MSE    & Mean Squared Error \\
PSNR   & Peak Signal-to-Noise Ratio \\
ReLU   & Rectified Linear Unit \\
RMSE   & Root Mean Squared Error \\
SE     & Squeeze-and-Excitation \\
SSIM   & Structural Similarity Index \\
\end{tabular}
% ===== end chapters/abbreviations.tex =====

\cleardoublepage

% ---------------------------------------------------------------- main matter
\pagenumbering{arabic}
\pagestyle{fancy}
\renewcommand{\chaptermark}[1]{\markboth{#1}{}}
\renewcommand{\sectionmark}[1]{\markright{\textbf{Chapter \thechapter}}}
% ===== begin chapters/ch1_introduction.tex =====
\chapter{Introduction}\label{ch:intro}

\section{Motivation}\label{sec:motivation}

An old photograph is rarely damaged in just one way. Take a print shot
through a soft lens on grainy film, scanned at low resolution and saved as a
heavily compressed JPEG. It carries all four defects, and each one hits a
different part of the spectrum. Optical softening removes high spatial
frequencies. Resampling throws away detail above a resolution limit and adds
interpolation artefacts. Sensor and film noise spreads energy roughly evenly
across frequencies. In natural images the true signal is weakest at high
frequencies, so that is where noise dominates. JPEG quantises the mid and
high frequency coefficients of each $8\times8$ block, which leaves blocking
and ringing.

Most learned restoration networks treat this mixture as a single problem.
DnCNN \cite{zhang2017dncnn}, a deep convolutional network, sees the whole
degraded image and learns to undo every corruption in it without being told
which is which. Transformer restorers such as Restormer
\cite{zamir2022restormer} work very well, but they have tens of millions of
parameters. An archive, a library or a hobbyist with modest hardware cannot
run them.

We started from a simple idea. If each corruption sits in its own frequency
band, we could split the image into bands explicitly and let a small network
pick a filter for every band at every location. A large kernel suits smooth,
slowly varying structure. A small kernel responds to fine detail and to the
high-frequency part of the noise. A learned gate chooses between them pixel
by pixel. We call the result \fcsg{}, a frequency-conditioned, spatially
gated network. We hold it to a strict budget of fewer than five million
parameters and fewer than ten GFLOPs for a $256\times256$ input.

\section{Problem statement}\label{sec:problem}

Let $f \in [0,1]^{H\times W\times3}$ be a clean RGB image and let
\begin{equation}
	x = \mathcal{D}(f;\, \theta), \qquad
	\theta = (\sigma_\text{blur},\, s,\, \sigma_n,\, q)
	\label{eq:degradation}
\end{equation}
be a degraded observation. $\mathcal{D}$ applies Gaussian blur of width
$\sigma_\text{blur}$, bicubic down- and up-sampling by a factor $s$, additive
Gaussian noise with standard deviation $\sigma_n$ and JPEG compression at
quality $q$. We draw $\theta$ at random for every image and never give it to
the restoration model. The task is to learn a mapping $g_\phi$ whose output
$\hat f = g_\phi(x)$ is close to $f$ in PSNR, SSIM and LPIPS. The network
must also meet two hard limits:
\begin{itemize}
	\item fewer than $5\times10^{6}$ trainable parameters, and
	\item fewer than 10 reported GFLOPs for one forward pass on a
	$256\times256$ RGB input.
\end{itemize}

\section{Hypothesis and objectives}\label{sec:objectives}

We test one hypothesis. \emph{Apply a fixed radial frequency decomposition,
then route each band per pixel to the top two of several shared
convolutional experts with different kernel sizes. The learned routing will
change with image content and with the degradation, the trained network will
depend on it for output quality, and the whole model will stay inside the
budget.} We wrote this as a claim to check, not as a promise that the network
beats anything.

The work has four objectives.
\begin{enumerate}
	\item[\textbf{O1}] Build \fcsg{} and check that it behaves as designed.
	The checks cover exact band reconstruction, the parameter and FLOP
	limits, memory, gradient flow and overfitting a single image
	(Chapters~\ref{ch:method} and~\ref{ch:setup}).
	\item[\textbf{O2}] Get baselines we can trust. We reproduce a published
	FFDNet result and train DnCNN on a Gaussian task with a known outcome
	(Chapter~\ref{ch:setup}).
	\item[\textbf{O3}] Train \fcsg{}, DnCNN and two FFDNet variants under one
	matched composite degradation protocol. Compare quality, size, compute and
	measured latency (Chapter~\ref{ch:results}).
	\item[\textbf{O4}] Describe what the router learned. Where does it send
	each band, how does that relate to the sampled degradation, and how much
	does the trained model rely on it (Chapter~\ref{ch:routing})?
\end{enumerate}

\section{Contributions}\label{sec:contributions}

This report makes four concrete contributions.
\begin{itemize}
	\item A restoration network that routes frequency bands spatially, with
	97,088 parameters and 9.39 reported GFLOPs at $256\times256$. We also give
	the tests showing the code matches the equations.
	\item A matched comparison. We train DnCNN, blind FFDNet, oracle FFDNet and
	\fcsg{} from scratch on the same synthetic composite degradation, with the
	same schedule, the same number of updates and the same evaluation
	protocol.
	\item A routing analysis. We relate per-pixel expert weights to the
	recorded degradation parameters and measure what happens when we override
	the gate at inference.
	\item A plain account of how it turned out. \fcsg{} reaches 27.24\dB{},
	against 27.83\dB{} for DnCNN and 28.19\dB{} for blind FFDNet, with 5.7
	times fewer parameters than DnCNN. It is not faster on a GPU. Its low FLOP
	count did not become low latency.
\end{itemize}
We make no claim of state-of-the-art quality. Frequency decomposition and
expert routing both have a long history (Chapter~\ref{ch:literature}). We
test one specific combination of them under a tight budget.

\section{Organisation of the report}\label{sec:organisation}

Chapter~\ref{ch:literature} reviews restoration networks, frequency-domain
methods and mixture-of-experts routing, and places \fcsg{} among them.
Chapter~\ref{ch:method} describes the degradation model and then the network
one module at a time. Chapter~\ref{ch:setup} covers the data pipeline,
baselines, training and evaluation protocols, and the verification runs.
Chapter~\ref{ch:results} reports restoration quality and computational cost.
Chapter~\ref{ch:routing} looks at the learned routing.
Chapter~\ref{ch:conclusion} sums up the findings and their limits and lists
the experiments we would run next. Appendix~\ref{app:repro} records what
anyone needs to reproduce the work, and Appendix~\ref{app:per-image} lists
results for every image.
% ===== end chapters/ch1_introduction.tex =====

\thispagestyle{empty}\cleardoublepage
% ===== begin chapters/ch2_literature.tex =====
\chapter{Background and related work}\label{ch:literature}

We place \fcsg{} against three bodies of work: convolutional restoration
networks, restoration in the frequency domain, and mixture-of-experts
routing. The papers below report scores on different tasks, datasets and
noise protocols, so we quote them for context only. Our result tables hold
only models we trained and evaluated ourselves under a single protocol.

\section{Image restoration with CNNs}\label{sec:lit-cnn}

\paragraph{DnCNN.} Zhang et al.\ \cite{zhang2017dncnn} found that a plain
stack of convolution, batch normalisation and ReLU learns denoising much more
easily when it predicts the noise residual instead of the clean image. The
network subtracts its predicted residual from the input to form the output.
Our RGB version has 17 layers of width 64 and 558,403 parameters. DnCNN takes
no noise-level input and knows nothing about frequency bands, so it shows
what a conventional network does on our data. The original paper also
covers blind denoising and other restoration tasks, not just one fixed noise
level.

\paragraph{FFDNet.} FFDNet \cite{zhang2018ffdnet} adds two ideas. It
rearranges the input into four half-resolution sub-images before any
convolution, which cuts computation by about four times. It also
concatenates a noise-level map to those sub-images, so one network covers a
range of noise strengths. The colour model has 12 layers of width 96. The
author's KAIR inference release \cite{kair} folds batch normalisation into the
convolutions and has 852,108 parameters. The paper reports 31.21\dB{} on
CBSD68 at $\sigma=25$ with unclipped Gaussian noise. FFDNet is our second
baseline. Its noise map describes Gaussian noise and nothing else. It carries
no information about blur or JPEG artefacts.

\paragraph{Larger restorers.} NAFNet \cite{chen2022nafnet} matches far more
complex designs with simplified blocks that drop conventional nonlinear
activations. New architecture alone, in other words, says little about
progress. Restormer \cite{zamir2022restormer} sits at the heavy end, with
transposed channel attention, gated feed-forward blocks and about 26 million
parameters in its deblurring setup. Both are far outside our budget. We cite
them as context, not as competitors.

\section{Frequency-domain restoration}\label{sec:lit-freq}

\paragraph{MWCNN.} Liu et al.\ \cite{liu2018mwcnn} swap the pooling and
upsampling of a U-Net for discrete wavelet and inverse wavelet transforms.
The transforms are fixed and have no parameters, so all learned capacity
lives in the convolutional backbone. MWCNN processes the subbands together,
and its ablations argue against treating each subband on its own. That
finding is why \fcsg{} has a cross-band fusion stage. MWCNN shows that an
explicit frequency split helps restoration. It does not send subbands to
different kernel sizes.

\paragraph{Fast Fourier Convolution.} FFC \cite{chi2020ffc} divides the
channels between a local convolutional branch and a global branch that works
in the Fourier domain. Every layer then has an image-wide receptive field.
FFC is an operator, not a restoration model, and its cost depends on the
backbone it goes into. \fcsg{} differs because it brings each masked band
back to the spatial domain before any learned layer touches it.

\paragraph{SFNet and FSNet.} Cui et al.\ wrote the closest prior work.
SFNet \cite{cui2023sfnet} selects frequency content with a content-dependent
local decomposition and channel attention. A second module pools globally and
in windows to widen the receptive field at low cost. FSNet \cite{cui2024fsnet}
applies the idea to six restoration tasks, with several receptive-field
scales in one U-shaped network. Both already pick frequencies locally and
adaptively. What we test is narrower. \fcsg{} uses \emph{fixed} radial masks
and makes its learned, spatial decision among shared $7\times7$, $5\times5$
and $3\times3$ experts. That decision can then be compared with known
degradation parameters. Fixed bands alone are not an advance over SFNet or
FSNet, and we do not claim they are.

\paragraph{FFTformer.} Kong et al.\ \cite{kong2023fftformer} move
attention-related operations into the frequency domain for deblurring. Their
GoPro model has 16.6 million parameters. They process Fourier blocks locally,
which shows that the spatial extent of the transform is a design choice in
its own right. The same issue comes up with our tiled inference
(Section~\ref{sec:latency}).

\paragraph{Spectral bias.} Rahaman et al.\ \cite{rahaman2019spectral} show
that ReLU networks fit the low-frequency parts of a target before the
high-frequency parts. That is a reason to inspect errors band by band. It
does not prove that a low band needs a large kernel, or that separate experts
will specialise. We treat both as open questions for the experiments.

\section{Mixture of experts and gating}\label{sec:lit-moe}

Shazeer et al.\ \cite{shazeer2017moe} introduced the sparsely gated
mixture-of-experts layer. A noisy top-$k$ gate picks a few experts for each
example. Extra importance and load penalties stop the gate from collapsing
onto a few favourites. These penalties balance usage over a batch. Maximising
the routing entropy of every example is a different thing, and it would push
each location toward uniform weights, which works against specialisation.
Fedus et al.\ \cite{fedus2022switch} cut routing to one expert per token in
the Switch Transformer and discuss capacity, load balance and numerical
stability at very large scale.

Two lessons from this work shape our design and limit our claims. First, the
\fcsg{} gate is spatial. It outputs a routing map over pixels, not one
decision per image. Second, the speed argument for sparse experts rests on
\emph{conditional computation}, where experts with zero weight never run. Our
code runs every expert and then mixes the outputs with sparse weights. The
sparsity lives in the weights and in the analysis. It does not make anything
faster, and every FLOP and latency figure we report includes all experts.

The fusion stage uses a squeeze-and-excitation block \cite{hu2018se}. It pools
features globally, squeezes them through a small bottleneck and predicts a
scale for each channel. In \fcsg{} it reweights the nine band channels before
the projection back to RGB, and it costs 96 parameters.

\section{Evaluation metrics}\label{sec:lit-metrics}

\paragraph{PSNR.} For pixel values in $[0,1]$, the peak signal-to-noise
ratio is
\begin{equation}
	\mathrm{PSNR}(\hat f, f) = 10\log_{10}\frac{1}{\mathrm{MSE}(\hat f, f)},
	\qquad
	\mathrm{MSE} = \frac{1}{3HW}\sum_{c,u,v}\big(\hat f_{c,u,v} - f_{c,u,v}\big)^2 .
	\label{eq:psnr}
\end{equation}
We compute PSNR per image over all three RGB channels and report the mean
over the dataset.

\paragraph{SSIM.} The structural similarity index \cite{wang2004ssim}
compares local means, variances and covariances inside a sliding window. We
use an $11\times11$ Gaussian window with $\sigma=1.5$ and constants
$C_1=0.01^2$ and $C_2=0.03^2$, then average over channels and positions.

\paragraph{LPIPS.} LPIPS \cite{zhang2018lpips} is a distance between
normalised deep features of two images. Its channel weights were fitted to
human judgements of perceptual similarity. We use the AlexNet variant with
inputs scaled to $[-1,1]$. Lower is better. LPIPS and PSNR can disagree,
because looking similar and matching pixel for pixel are different things.
The LPIPS network is only a measuring tool, and we do not count its
parameters toward any model.

\section{Research gap}\label{sec:gap}

Table~\ref{tab:gap} sets \fcsg{} beside the methods above. It combines fixed
frequency bands, shared experts that differ in kernel size, per-pixel top-2
routing and a budget under 0.1 million parameters. None of the methods we
reviewed sits at exactly this point. Each part, though, already exists. Our
contribution is an empirical study of how the parts behave together, not a
new building block.

\begin{table}[htbp]
	\centering
	\caption{How \fcsg{} compares with closely related methods.}
	\small
	\label{tab:gap}
	\begin{tabular}{@{}lcccl@{}}
		\toprule
		Method & Freq.\ split & Learned selection & Spatial & Notes \\
		\midrule
		DnCNN \cite{zhang2017dncnn}      & --         & --              & --        & residual CNN \\
		FFDNet \cite{zhang2018ffdnet}    & --         & --              & --        & noise-map input \\
		MWCNN \cite{liu2018mwcnn}        & wavelet    & --              & --        & joint subbands \\
		SFNet \cite{cui2023sfnet}        & learned    & frequency       & local     & closest prior work \\
		Sparse MoE \cite{shazeer2017moe} & --         & experts         & --        & per-example routing \\
		\fcsg{} (ours)                   & radial FFT & experts, top-2  & per pixel & 97k parameters \\
		\bottomrule
	\end{tabular}
\end{table}
% ===== end chapters/ch2_literature.tex =====

\thispagestyle{empty}\cleardoublepage
% ===== begin chapters/ch3_method.tex =====
\chapter{Proposed method: \fcsg}\label{ch:method}

We first describe the synthetic degradation used for training and
evaluation, then walk through the network one module at a time. The notation
matches the code in \texttt{models/fcsg.py} and \texttt{models/blocks.py}.

\section{Degradation model}\label{sec:degradation}

An archival photograph picks up damage in a rough physical order. The lens
softens it, scanning or resizing loses resolution, the sensor or film adds
noise, and lossy storage compresses what remains. Our pipeline applies four
matching operations in the same order (Table~\ref{tab:degradation}).
\begin{table}[htbp]
	\centering
	\caption{Composite degradation pipeline. We sample the parameters independently for every image.}
	\label{tab:degradation}
	\begin{tabular}{@{}lll@{}}
		\toprule
		Stage & Parameter & Distribution \\
		\midrule
		Gaussian blur        & $\sigma_\text{blur}$ & $\mathcal{U}(0.5,\, 2.0)$ pixels \\
		Down/up resampling   & scale $s$            & $\mathcal{U}(1.0,\, 2.0)$, bicubic both ways \\
		Additive noise       & $\sigma_n$           & $\mathcal{U}(5,\, 50)/255$, clipped to $[0,1]$ \\
		JPEG compression     & quality $q$          & integer, uniform on $\{30,\dots,95\}$, 4:4:4 \\
		\bottomrule
	\end{tabular}
\end{table}

Pillow's Gaussian filter does the blur. Resampling shrinks the image by $s$
with bicubic interpolation and then enlarges it back, so input and output
always share a size. The noise is additive white Gaussian noise in 8-bit
units, clipped to the valid range. Pillow encodes the JPEG with chroma
subsampling turned off. Before the blur and again before JPEG encoding, we
quantise the image to 8 bits, as a real file would be.

Each sample stores its four parameters next to the degraded and clean
images. The network never sees them. We keep them so that after training we
can ask whether the router reacts to the degradation
(Chapter~\ref{ch:routing}). Noise runs from 5 to 50 grey levels and JPEG
quality goes as low as 30, and the mean input PSNR on our validation set
comes out at 20.99\dB{}.

\section{Architecture overview}\label{sec:arch}

Figure~\ref{fig:architecture} shows the data flow. The network splits the
degraded image $x$ into a low, a mid and a high frequency band, and each band
is an ordinary RGB image. Three shared experts with $7\times7$, $5\times5$
and $3\times3$ kernels process every band. At each pixel and for each band, a
spatial gate gives every expert a weight and keeps the two largest. The
weighted expert outputs add up to one restored image per band:
\begin{equation}
	y_b(u,v) = \sum_{e\in\{7,5,3\}} w_{b,e}(u,v)\; E_e(x_b)(u,v),
	\qquad b\in\{L,M,H\}.
	\label{eq:mix}
\end{equation}
The network concatenates the three band outputs into nine channels,
reweights them with a squeeze-and-excitation block, projects them back to RGB
and refines the result. It predicts a residual and adds it to the input:
\begin{equation}
	\hat{f} = x + \mathrm{Refine}\big(\mathrm{SEFusion}([\,y_L, y_M, y_H\,])\big).
	\label{eq:out}
\end{equation}

\begin{figure}[htbp]
	\centering
	\includegraphics[width=\linewidth]{architecture.pdf}
	\caption{\fcsg{} data flow. The three bands share expert parameters. Before
	each expert call, the band dimension is folded into the batch dimension.}
	\label{fig:architecture}
\end{figure}

\section{Frequency decomposition}\label{sec:freqdecomp}

We split the image in the Fourier domain with fixed masks, then bring each
masked band straight back to the spatial domain. We chose this early, and
for a reason. An earlier draft of the proposal fed masked complex spectra
directly into the convolutional experts, and that fails in two ways. A
convolution over a spectrum has no clean meaning, because multiplication in
one domain is convolution in the other. The argument that large kernels suit
smooth structure stops making sense there. Complex inputs also cause
trouble. Split into two real channels, the phase is discontinuous and
training becomes unstable. Kept complex, they need tensor support that many
PyTorch operations lack. Masking in frequency and filtering in space keeps
the split exact and lets every expert be a plain real-valued CNN.

Let $X = \mathcal{F}(x)$ be the two-dimensional real FFT of the image with
orthonormal scaling, and let $r = \sqrt{f_x^2 + f_y^2}$ be the radial
frequency in cycles per pixel. We build the frequency grid with
\texttt{fftfreq} and \texttt{rfftfreq} in the unshifted order the FFT uses.
Each axis reaches the Nyquist limit at $0.5$ cycles per pixel. A soft
low-pass mask with cutoff $c$ and transition width $\Delta = 0.05$ is a
raised cosine:
\begin{equation}
	t_c(r) = \mathrm{clip}\!\left(\frac{r-c}{\Delta} + \tfrac12,\ 0,\ 1\right),
	\qquad
	\mathrm{LP}_c(r) = \tfrac12\big(1 + \cos(\pi\, t_c(r))\big).
	\label{eq:lowpass}
\end{equation}
With cutoffs $r_1 = 0.10$ and $r_2 = 0.25$, the three band masks are
\begin{equation}
	M_L = \mathrm{LP}_{r_1}, \qquad
	M_M = \mathrm{LP}_{r_2} - \mathrm{LP}_{r_1}, \qquad
	M_H = 1 - \mathrm{LP}_{r_2},
	\label{eq:masks}
\end{equation}
and each band image is $x_b = \mathcal{F}^{-1}(M_b \cdot X)$.

Three properties of this split matter. The first is exactness. The masks sum
to one at every frequency and the FFT is linear, so $x_L + x_M + x_H = x$.
The split loses nothing and adds nothing. In single precision we measured a
maximum reconstruction error of $6.0\times10^{-7}$. The second is that the
cutoffs are in cycles per pixel, not FFT bin indices. A given bin index means
one physical frequency at $128\times128$ and another at full resolution, so
cutoffs set in bins would drift between training crops and test images. With
normalised radii, a sinusoid lands in the same band at both sizes. We checked
this at 128 and 256 pixels. The third is that the transitions are smooth. A
hard step in frequency causes ringing in space, and the network would waste
capacity cleaning up after its own preprocessing.

The FFT stays in single precision even under mixed-precision training,
because half-precision CUDA FFTs only handle power-of-two sizes.
Figure~\ref{fig:masks} plots the three masks.

\begin{figure}[htbp]
	\centering
	\includegraphics[width=0.9\linewidth]{frequency_masks.pdf}
	\caption{Raised-cosine band masks with cutoffs at 0.10 and 0.25 cycles per
	pixel and a transition width of 0.05. Right, the masks drawn on a centred
	frequency plane. The computation itself uses unshifted order.}
	\label{fig:masks}
\end{figure}

\section{Shared kernel-size experts}\label{sec:experts}

The three experts, $E_7$, $E_5$ and $E_3$, differ only in kernel size. Each
one processes all three bands. The bands therefore share parameters, and
adding bands would not add experts.

An expert starts with a lossless $2\times2$ pixel unshuffle
\cite{shi2016espcn}, as in FFDNet. This turns a three-channel image into
twelve channels at half resolution. A $1\times1$ convolution widens these to
48 channels, followed by a ReLU. Next come four residual blocks. Each block
has two depthwise-separable stages \cite{chollet2017xception}, each made of a
$k\times k$ depthwise convolution and a $1\times1$ pointwise convolution, with
a ReLU after the first stage. A $1\times1$ tail maps back to twelve channels
and adds the unshuffled input. A pixel shuffle then restores full resolution.
For odd image sizes the expert pads by replication and crops afterwards.

The budget depends on two choices, separable convolutions and
half-resolution processing. Each image goes through nine expert passes, one
per band for each of three experts. Dense full-resolution convolutions at
width 48 would blow well past ten GFLOPs.

We name the experts by kernel size and nothing else. Calling one a ``texture
expert'' or an ``edge expert'' would be a claim about what it learned, and no
attribution experiment of ours supports such a claim.

\section{Spatial top-2 gate}\label{sec:gate}

A gate that decided once per image would output nine numbers from a few
global statistics. That is closer to a lookup table than a router, and it
would probably settle on a near-constant choice. The \fcsg{} gate decides per
location instead, so a smooth sky and a textured roof in the same photograph
can get different treatment.

The gate stacks the three RGB bands into nine channels. In non-overlapping
$16\times16$ windows it computes the local mean and the local mean of squares
of every channel. A two-layer $1\times1$ network ($18\to24\to9$, with a ReLU
between the layers) maps these eighteen statistics to nine logits $a_{b,e}$,
one per band and expert. The gate divides the logits by a temperature
$\tau = 1$ and upsamples them bilinearly to full resolution. At each pixel
and for each band it keeps the two largest logits and normalises them with a
softmax. The third expert gets a weight of zero:
\begin{equation}
	w_{b,e}(u,v) =
	\begin{cases}
		\dfrac{\exp a_{b,e}(u,v)}{\sum_{j\in\mathcal{T}_b(u,v)} \exp a_{b,j}(u,v)}
		& e \in \mathcal{T}_b(u,v),\\[2ex]
		0 & \text{otherwise},
	\end{cases}
	\label{eq:topk}
\end{equation}
where $\mathcal{T}_b(u,v)$ holds the two highest-scoring experts. The gate
also returns $p_{b,e}$, the dense softmax over all three experts, which feeds
the entropy term of the loss. The whole gate has 681 parameters.
Figure~\ref{fig:gate} shows the expert and gate operations.

\begin{figure}[htbp]
	\centering
	\includegraphics[width=0.9\linewidth]{expert_gate.pdf}
	\caption{Expert and gate operations. Experts use pixel rearrangement and
	depthwise-separable residual blocks. The gate pools band statistics in
	$16\times16$ windows.}
	\label{fig:gate}
\end{figure}

As Section~\ref{sec:lit-moe} pointed out, every expert runs before mixing.
Where an expert's weight is zero, the network still computes its output and
then multiplies it by zero. Saving that work would take tiled conditional
dispatch, which we did not build.

\section{Fusion and residual refinement}\label{sec:fusion}

The network concatenates the three weighted band outputs into nine channels.
A squeeze-and-excitation block \cite{hu2018se} pools them globally, passes
them through a $9\to3\to9$ bottleneck with a sigmoid, and scales each
channel. A $1\times1$ convolution projects the result to three channels. Three
$3\times3$ convolutions ($3\to16\to16\to3$, with ReLUs between them) refine
it, and the output is added to the degraded input as in
Equation~\eqref{eq:out}. Predicting a residual instead of the image is the
main lesson of DnCNN \cite{zhang2017dncnn}. It also means the network starts
out close to the identity.

\section{Training objective}\label{sec:loss}

The loss has three terms:
\begin{equation}
	\mathcal{L} = \mathcal{L}_\text{char}
	+ 0.05\,\mathcal{L}_\text{freq}
	+ \lambda_\text{ent}(t)\,\mathcal{L}_\text{ent}.
	\label{eq:loss}
\end{equation}

The Charbonnier loss \cite{charbonnier1994} is a smooth form of L1,
\begin{equation}
	\mathcal{L}_\text{char} = \operatorname{mean}\sqrt{(\hat f - f)^2 + \varepsilon^2},
	\qquad \varepsilon = 10^{-3}.
\end{equation}
Near zero it acts like L2, and for large errors it acts like L1.

The frequency term is the mean absolute value of the orthonormal FFT of the
error, $\mathcal{L}_\text{freq} = \operatorname{mean}\lvert
\mathcal{F}(\hat f - f)\rvert$. It gives every frequency the same weight.
Pixel losses tend to favour the dominant low frequencies, and this term
pushes back against that.

The entropy term is the negative dense routing entropy, summed over bands and
averaged over the batch and over pixels:
\begin{equation}
	\mathcal{L}_\text{ent} = \sum_b \operatorname{mean}_{n,u,v}\sum_e p_{b,e}\log p_{b,e}.
\end{equation}
Minimising this term maximises entropy. Early in training that keeps all
three experts in use, so none of them goes short of gradient. A fixed entropy
bonus, though, would keep pulling the router toward uniform weights for the
whole run, which is the opposite of what we want to see. So the coefficient
decays linearly, $\lambda_\text{ent}(t) = 0.01\,(1 - t/T)$, and hits zero at
the last update $T$. The sign of this term is easy to get wrong, so the
verification suite checks that it falls as the weights approach uniform.
DnCNN and FFDNet train with the same Charbonnier and frequency terms and no
entropy term.

\section{Model size and complexity}\label{sec:size}

Table~\ref{tab:params} breaks down the 97,088 parameters. The three experts
hold nearly all of them. The gate, fusion and refinement stages together come
to under four thousand.

\begin{table}[htbp]
	\centering
	\caption{\fcsg{} parameters by module.}
	\label{tab:params}
	\begin{tabular}{@{}lr@{}}
		\toprule
		Module & Parameters \\
		\midrule
		Shared $7\times7$ expert            & 39,228 \\
		Shared $5\times5$ expert            & 30,012 \\
		Shared $3\times3$ expert            & 23,868 \\
		Spatial gate ($18\to24\to9$)        & 681 \\
		SE fusion and $9\to3$ projection    & 96 \\
		Residual refinement                 & 3,203 \\
		\midrule
		Total                               & 97,088 \\
		\bottomrule
	\end{tabular}
\end{table}

To count compute we profile multiply-accumulate operations (MACs) with
\texttt{thop} and count each MAC as two floating-point operations. For the
forward and inverse FFTs we add an estimate of $12 \times 5N\log_2 N$, with
$N = 256^2$. We leave out small elementwise operations. By this count a
$256\times256$ forward pass costs 4.66 GMACs, or 9.39 GFLOPs with the FFT
estimate, which is inside the 10 GFLOP limit. The baselines are counted the
same way. Chapter~\ref{ch:results} shows that this number turned out to
predict running time badly.
% ===== end chapters/ch3_method.tex =====

\thispagestyle{empty}\cleardoublepage
% ===== begin chapters/ch4_setup.tex =====
\chapter{Experimental setup and verification}\label{ch:setup}

The numbers in Chapter~\ref{ch:results} can be no better than the code and
protocol that produced them. Here we describe the data, the baselines, and
how we trained and evaluated every model. We then cover the checks we ran
before spending GPU hours on long training runs.

\section{Dataset and data pipeline}\label{sec:data}

Every experiment uses DIV2K \cite{agustsson2017div2k}, which has 800
high-resolution training images and 100 validation images. Training draws
only on the training split. Every result we report comes from the 100
validation images. We kept no separate test set, so these are validation
numbers. Chapter~\ref{ch:conclusion} comes back to this limitation.

Our first data loader decoded a full-resolution PNG for every crop. Training
crawled along at about 1.8 steps per second, with the GPU waiting on the CPU.
We replaced the loader with a cache of 32 fixed $128\times128$ crops per
training image, 25,600 tiles in total. They sit in one 1.26\,GB array of
8-bit pixels that we memory-map during training. Random flips and 90-degree
rotations augment the cached crops. Four loader workers still apply the
degradation of Section~\ref{sec:degradation} on the fly, so each epoch gets
freshly sampled blur, scale, noise and JPEG settings. The clean crops are a
fixed set, but the supply of degraded pairs never runs out
(Figure~\ref{fig:pipeline}).

\begin{figure}[htbp]
	\centering
	\includegraphics[width=\linewidth]{pipeline.pdf}
	\caption{Training data pipeline. For each augmented clean crop, a loader
	worker samples degradation parameters and returns the degraded input, the
	clean target and the parameters.}
	\label{fig:pipeline}
\end{figure}

Degrading on the fly helps training, but it makes exact inputs hard to
reproduce. So we also exported a fixed set of 5,000 degraded and clean pairs
generated with seed 1234, plus a manifest of their parameters. The
degradation code regenerates this set byte for byte. One CPU worker produces
about 233 samples per second.

\section{Baselines}\label{sec:baselines}

\paragraph{DnCNN.} Our RGB DnCNN has 17 convolutional layers of width 64
with batch normalisation. It predicts the residual and has 558,403
parameters.

\paragraph{FFDNet.} Our FFDNet copies the author's colour architecture
\cite{kair}, with a pixel unshuffle, 12 convolutional layers of width 96 and a
noise-level map concatenated to the input. It has 852,108 parameters. On our
composite task the noise map causes a problem, since Gaussian noise is only
one of four corruptions. We train two variants to deal with this.
\emph{Oracle} FFDNet gets the true sampled noise level $\sigma_n$ as its map.
No real user would know this value, so the oracle serves as an upper
reference and not a fair competitor. \emph{Blind} FFDNet gets a constant map
of $27.5/255$, the middle of the noise range, in both training and
evaluation. It sees exactly what DnCNN and \fcsg{} see.

We train all four models from scratch on the same composite degradation. We
do not compare against numbers published for other tasks or datasets.

\section{Training protocol}\label{sec:training}

All models share one configuration (Table~\ref{tab:hparams}). Only the model
changes, along with the noise map for FFDNet. We had planned 300,000 updates
and cut this to 200,000 to fit our weekly GPU quota.

\begin{table}[htbp]
	\centering
	\caption{Training configuration, the same for every method.}
	\label{tab:hparams}
	\small
	\begin{tabular}{@{}ll@{}}
		\toprule
		Setting & Value \\
		\midrule
		Optimiser            & AdamW \cite{loshchilov2019adamw}, weight decay 0 \\
		Learning rate        & $2\times10^{-4}$, cosine decay \cite{loshchilov2017sgdr} to $10^{-6}$ \\
		Updates              & 200,000 \\
		Batch                & 16 crops of $128\times128$ RGB \\
		Loss                 & $\mathcal{L}_\text{char} + 0.05\,\mathcal{L}_\text{freq}$, plus entropy term for \fcsg{} \\
		Entropy coefficient  & 0.01, linear decay to 0 (\fcsg{} only) \\
		Precision            & FP16 mixed precision \cite{micikevicius2018amp} with gradient scaling \\
		Validation           & 16 fixed crops every 2,000 updates \\
		Checkpoints          & every 2,000 updates (model, optimiser, scaler, step) \\
		Hardware             & NVIDIA Tesla T4 (Kaggle) \\
		\bottomrule
	\end{tabular}
\end{table}

We trained in Kaggle notebook sessions, which end after twelve hours, so each
run picks up from its newest checkpoint when it starts. A checkpoint holds the
model, optimiser, gradient scaler, step and configuration. It does not hold
the random number generator or sampler state, which means a resumed run is
not bit-identical to an unbroken one. We did not fix training seeds, and we
trained each model once. The training loop logged 10.06 hours for \fcsg{},
6.66 for DnCNN, 2.04 for blind FFDNet and 2.33 for oracle FFDNet.
Figure~\ref{fig:schedules} shows the learning-rate and entropy schedules.

\begin{figure}[htbp]
	\centering
	\includegraphics[width=0.85\linewidth]{schedules.pdf}
	\caption{Learning-rate and entropy-coefficient schedules over 200,000 updates.}
	\label{fig:schedules}
\end{figure}

\section{Evaluation protocol}\label{sec:evalproto}

Final results cover all 100 DIV2K validation images. We degrade each image
once, using its index as the seed, so every model gets identical degraded
inputs. Inference runs in tiles of $256\times256$ with 32 pixels of overlap
(stride 224). The last tile in each direction lines up with the image border,
and overlapping outputs are averaged with equal weight. Tiling keeps memory
low on large images, and Section~\ref{sec:latency} measures what it costs in
accuracy. We clip restored values to $[0,1]$ and then compute PSNR, SSIM and
LPIPS as in Section~\ref{sec:lit-metrics}.

During training we tracked something else, the PSNR on 16 fixed
$128\times128$ validation crops. That number cannot be compared with the
full-image result. For \fcsg{}, the final crop PSNR is 30.59\dB{} and the
full-image PSNR is 27.24\dB{}. We use crop numbers only to plot convergence.

\section{Implementation verification}\label{sec:verification}

Before any long run, we ran a fixed list of checks in one notebook on the
same Kaggle T4 (Python 3.12, PyTorch 2.10 \cite{paszke2019pytorch}, CUDA 12.8).
The notebook halts at the first failure. Table~\ref{tab:gates} lists each
check and the value we measured.

\begin{table}[htbp]
	\centering
	\caption{Checks run before training (Tesla T4, run of 25 September 2026).}
	\small
	\label{tab:gates}
	\begin{tabular}{@{}lll@{}}
		\toprule
		Check & Criterion & Measured \\
		\midrule
		FFDNet weights, CBSD68 $\sigma{=}25$ & within 0.5\dB{} of 31.21 & 31.220\dB \\
		Band reconstruction & $x_L{+}x_M{+}x_H = x$ & max error $6.0\times10^{-7}$ \\
		Frequency units & same band at 128 and 256 px & passed \\
		Parameter budget & $< 5$M & 97,088 \\
		Compute budget & $< 10$ GFLOPs at $256^2$ & 9.39 \\
		Memory, forward + backward & $< 4$ GB & 0.60 GB \\
		Shapes, gradients, odd sizes (AMP) & no errors & passed \\
		One-image overfit, 200 steps & MSE $<0.01$, $<\frac14$ initial & ratio 0.18 \\
		Checkpoint restore & identical outputs & passed \\
		Degradation reproducibility & byte-identical & passed \\
		\bottomrule
	\end{tabular}
\end{table}

Two checks need a comment. The band reconstruction check is one line of
code, and it still catches almost every sign, shift and normalisation bug an
FFT implementation can have. The overfit check asks the full model to
memorise a single image in 200 steps. If it cannot, there is a bug, and
learning that in a few seconds beats learning it after ten hours of training.
\fcsg{} brought the MSE down from 0.0261 to 0.0046 (Figure~\ref{fig:overfit}).
We measured memory at batch size 1, so that figure says nothing about peak
memory at the training batch size.

\begin{figure}[htbp]
	\centering
	\includegraphics[width=0.75\linewidth]{overfit.pdf}
	\caption{Single-image overfit test. MSE over 200 updates.}
	\label{fig:overfit}
\end{figure}

\section{Baseline reproduction}\label{sec:baseline-repro}

\paragraph{FFDNet.} To check our FFDNet code and evaluation path against a
published number, we loaded the author's pretrained colour weights
\cite{kair}. We ran them on all 68 CBSD68 images \cite{martin2001bsd} with
unclipped Gaussian noise at $\sigma=25$, as the author did. We got
31.220\dB{}. The paper \cite{zhang2018ffdnet} reports 31.21\dB{}, a gap of
0.01\dB{}. This confirms the architecture, the weight loading and the
evaluation code. It does not reproduce FFDNet's training, which is why we
train every composite model from scratch.

\paragraph{DnCNN.} To test our training loop end to end, we trained DnCNN on
clipped Gaussian noise at $\sigma=25$ before starting the composite task. At
checkpoint 286,000 it scored 32.65\dB{} on DIV2K validation and 31.04\dB{} on
CBSD68 (Table~\ref{tab:gaussian}). These numbers use clipped noise, so they
do not line up directly with the unclipped protocol of the FFDNet paper. They
show that our training and evaluation code produce results of the expected
size on a well-known problem.

\begin{table}[htbp]
	\centering
	\caption{DnCNN on clipped Gaussian noise, $\sigma=25$, checkpoint 286,000.}
	\label{tab:gaussian}
	\begin{tabular}{@{}lrrr@{}}
		\toprule
		Dataset & Images & Input PSNR & Output PSNR \\
		\midrule
		CBSD68           & 68  & 20.53\dB & 31.04\dB \\
		DIV2K validation & 100 & 20.70\dB & 32.65\dB \\
		\bottomrule
	\end{tabular}
\end{table}
% ===== end chapters/ch4_setup.tex =====

\thispagestyle{empty}\cleardoublepage
% ===== begin chapters/ch5_results.tex =====
\chapter{Restoration results and computational cost}\label{ch:results}

This chapter covers how well the four trained models restore the 100 DIV2K
validation images, how quality relates to size and compute, and what each
model costs to run. We trained every model here once, for 200,000 updates,
under the protocol of Chapter~\ref{ch:setup}.

\section{Quantitative comparison}\label{sec:main-table}

Table~\ref{tab:main} holds the main result. Every model starts from the same
degraded inputs, whose mean PSNR is 20.99\dB{}.

\begin{table}[htbp]
	\centering
	\caption{Composite restoration on 100 DIV2K validation images after 200,000
	updates. Gain is output PSNR minus input PSNR (input 20.992\dB{}). GFLOPs
	at $256\times256$.}
	\label{tab:main}
	\small
	\begin{tabular}{@{}lrrrrrr@{}}
		\toprule
		Method & PSNR$\uparrow$ & Gain & SSIM$\uparrow$ & LPIPS$\downarrow$ & Params & GFLOPs \\
		\midrule
		DnCNN                     & 27.831 & 6.838 & 0.7795 & 0.3242 & 558,403 & 73.43 \\
		FFDNet, blind             & 28.190 & 7.197 & 0.7899 & 0.3176 & 852,108 & 27.89 \\
		FFDNet, oracle$^\dagger$  & 28.240 & 7.247 & 0.7920 & 0.3162 & 852,108 & 27.89 \\
		\fcsg{}                   & 27.243 & 6.250 & 0.7529 & 0.3961 &  97,088 &  9.39 \\
		\bottomrule
		\multicolumn{7}{@{}l}{\footnotesize $^\dagger$Receives the true sampled noise level. An upper reference, not a fair competitor.}
	\end{tabular}
\end{table}

\fcsg{} lifts the input by 6.25\dB{}, so it does learn real restoration. It
is still the weakest of the four on all three metrics, 0.59\dB{} behind DnCNN
and 0.95\dB{} behind blind FFDNet. The LPIPS gap is wider than the PSNR gap.
\fcsg{} scores 0.396 against 0.318 for blind FFDNet, which hints that it
leaves more visible residue than the PSNR shortfall alone would predict.

Size and compute change the picture (Table~\ref{tab:tradeoff}). \fcsg{} comes
within 0.6\dB{} of DnCNN with 82.6\% fewer parameters and 87.2\% fewer
reported FLOPs. It comes within 1\dB{} of both FFDNet variants with 88.6\%
fewer parameters and 66.3\% fewer FLOPs.

\begin{table}[htbp]
	\centering
	\caption{\fcsg{} measured against each baseline.}
	\label{tab:tradeoff}
	\begin{tabular}{@{}lrrr@{}}
		\toprule
		Compared with & PSNR difference & Fewer parameters & Fewer GFLOPs \\
		\midrule
		DnCNN           & $-0.588$\dB & 82.6\% & 87.2\% \\
		FFDNet, blind   & $-0.947$\dB & 88.6\% & 66.3\% \\
		FFDNet, oracle  & $-0.997$\dB & 88.6\% & 66.3\% \\
		\bottomrule
	\end{tabular}
\end{table}

The baselines themselves tell us two things. Blind FFDNet is the strongest
fair model, even though its noise map is a constant with no information in
it. Its lead over DnCNN must come from the architecture. Our best guess is
the half-resolution processing, which widens the effective receptive field of
every layer. Noise conditioning cannot explain it. The oracle map, for its
part, adds only 0.05\dB{} over the blind constant. Knowing the noise level is
worth little once blur, resampling and JPEG artefacts are mixed in, because
the map describes one corruption out of four.

Figure~\ref{fig:benchmark} plots the three metrics side by side.

\begin{figure}[htbp]
	\centering
	\includegraphics[width=\linewidth]{benchmark.pdf}
	\caption{Mean PSNR, SSIM and LPIPS over the 100 validation images. Each bar
	is a single training run.}
	\label{fig:benchmark}
\end{figure}

We trained each model once without fixing the seed, so we cannot tell the
effect of architecture apart from run-to-run noise. A difference of a few
hundredths of a decibel, like the one between the two FFDNet variants, means
little. The gap between \fcsg{} and the baselines is a different matter. It
is large and it holds across images (Section~\ref{sec:per-image}), and we do
not expect a second seed to reverse it.

\section{Per-image analysis}\label{sec:per-image}

A good average can hide a model that wins most images and fails badly on a
few. Table~\ref{tab:wins} compares \fcsg{} with each baseline one image at a
time.

\begin{table}[htbp]
	\centering
	\caption{Per-image PSNR of \fcsg{} against each baseline, 100 images.}
	\label{tab:wins}
	\begin{tabular}{@{}lccc@{}}
		\toprule
		Compared with & Wins / ties / losses & Mean difference & Range \\
		\midrule
		DnCNN          & 6 / 1 / 93  & $-0.59$\dB & $-2.75$ to $+2.07$\dB \\
		FFDNet, blind  & 0 / 0 / 100 & $-0.95$\dB & $-3.12$ to $-0.04$\dB \\
		FFDNet, oracle & 0 / 0 / 100 & $-1.00$\dB & $-3.48$ to $-0.07$\dB \\
		\bottomrule
	\end{tabular}
\end{table}

The shortfall is systematic. Both FFDNet variants beat \fcsg{} on all 100
images. Against DnCNN, \fcsg{} wins six images, one of them by 2.07\dB{}, so
the two convolutional designs break down in different places. Take image
0801. DnCNN barely moves it, from 28.09 to 28.22\dB{}, while \fcsg{} reaches
29.32\dB{}. DnCNN is also the only model that makes some images worse. Three
of its outputs score below their inputs.

Image 0828 is the hardest for every model. Its degraded input scores just
17.10\dB{}. \fcsg{} reaches 17.93\dB{} on it, and the baselines land between
18.61 and 18.83\dB{}. Our setup cannot say which of the four sampled
corruptions makes this image so hard. Figure~\ref{fig:per-image} shows the
paired differences and the error distributions, and
Appendix~\ref{app:per-image} gives the full table.

\begin{figure}[htbp]
	\centering
	\includegraphics[width=\linewidth]{per_image_errors.pdf}
	\caption{Left, per-image PSNR difference between \fcsg{} and each baseline.
	Right, sorted per-image RMSE for every model, in 8-bit units.}
	\label{fig:per-image}
\end{figure}

\section{Qualitative results}\label{sec:qualitative}

Figures~\ref{fig:qual-a} and~\ref{fig:qual-b} show the top-left
$320\times320$ region of validation images 0801 and 0802 for each model.
All models get the identical degraded input. The PSNR labels refer to the
whole image, not to the region shown.

\begin{figure}[htbp]
	\centering
	\begin{subfigure}{\linewidth}
		\centering
		\includegraphics[width=0.92\linewidth]{qualitative_fcsg.pdf}
		\caption{\fcsg{}}
	\end{subfigure}\\[1ex]
	\begin{subfigure}{\linewidth}
		\centering
		\includegraphics[width=0.92\linewidth]{qualitative_dncnn.pdf}
		\caption{DnCNN}
	\end{subfigure}
	\caption{Degraded input, restoration and clean target for \fcsg{} and DnCNN.}
	\label{fig:qual-a}
\end{figure}

\begin{figure}[htbp]
	\centering
	\begin{subfigure}{\linewidth}
		\centering
		\includegraphics[width=0.92\linewidth]{qualitative_blind.pdf}
		\caption{FFDNet, blind ($\sigma$ map fixed at 27.5/255)}
	\end{subfigure}\\[1ex]
	\begin{subfigure}{\linewidth}
		\centering
		\includegraphics[width=0.92\linewidth]{qualitative_oracle.pdf}
		\caption{FFDNet, oracle (true $\sigma$ map)}
	\end{subfigure}
	\caption{Degraded input, restoration and clean target for the two FFDNet variants.}
	\label{fig:qual-b}
\end{figure}

All four models clear most of the noise and blocking. None brings back the
fine texture that blur and downsampling destroyed. We expected that. The
detail is gone from the input, and a model trained on pixel losses predicts a
smooth average instead of making texture up. At this scale the models look
much alike, and the saved panels have no difference images. We therefore lean
on the numbers above, not on visual judgement.

\section{Training convergence}\label{sec:convergence}

Figure~\ref{fig:convergence} plots validation PSNR on the 16 fixed crops
against updates and against wall-clock time. By the end of training all four
curves have nearly levelled off. \fcsg{} finishes at 30.59\dB{} on the crops,
below DnCNN (31.66), blind FFDNet (31.97) and oracle FFDNet (32.05). The
order matches the full evaluation. With the curves this flat, nothing
suggests that more training under the same settings would close the gap.

\begin{figure}[htbp]
	\centering
	\includegraphics[width=\linewidth]{convergence.pdf}
	\caption{Validation PSNR on 16 fixed crops against updates and training
	time. Crop PSNR cannot be compared with full-image results.}
	\label{fig:convergence}
\end{figure}

The wall-clock panel shows a cost that update counts hide. \fcsg{} took 10.06
hours for 200,000 updates. DnCNN took 6.66 and FFDNet about two. We also ran
a separate throughput test without warm-up, validation or checkpointing. It
measured 5.84 steps per second for \fcsg{}, 9.41 for DnCNN and 26.67 for
FFDNet. Equal update counts did not mean equal GPU time, and \fcsg{} got the
most.

\section{Latency, memory and tiling}\label{sec:latency}

\paragraph{Latency.} We built a small model so it would be cheap to run.
Table~\ref{tab:latency} shows that on our GPU it is not. Each figure is the
median wall-clock time of one FP32 forward pass at batch size 1 on a Tesla
T4, with TF32 off, taken over 50 timed runs after 20 warm-up passes.

\begin{table}[htbp]
	\centering
	\caption{Median forward latency (ms) by input size, and peak inference memory at $256\times256$ (Tesla T4, FP32, batch 1).}
	\label{tab:latency}
	\begin{tabular}{@{}lrrrr@{}}
		\toprule
		Method & $128^2$ & $256^2$ & $512^2$ & Peak MB at $256^2$ \\
		\midrule
		\fcsg{}          & 5.10 & 19.40 & 75.59 & 51.3 \\
		DnCNN            & 4.22 & 18.04 & 72.47 & 37.8 \\
		FFDNet, blind    & 1.69 &  6.30 & 24.12 & 20.2 \\
		FFDNet, oracle   & 1.69 &  6.35 & 24.62 & 20.2 \\
		\bottomrule
	\end{tabular}
\end{table}

\fcsg{} is the slowest model at every input size. It trails DnCNN slightly
and runs three times slower than FFDNet, although it needs an eighth of
DnCNN's reported FLOPs and a third of FFDNet's. It also uses the most memory.
We did not profile individual operators, so what follows is our reading of
the architecture, not a measurement. Every forward pass runs nine full expert
passes, because experts with zero weight are not skipped. Depthwise
convolutions do very little arithmetic for each byte they read, so on a GPU
memory traffic sets their speed, not the FLOP count. The network is built
from many small layers, which makes kernel launch overhead a large share of
the total. The FFT and inverse FFT add work that our FLOP estimate counts only
roughly. FFDNet runs a single stack of dense $3\times3$ convolutions at half
resolution, which is exactly the workload GPUs are designed for.

The FLOP budget we fixed at the start was the wrong stand-in for speed on
this hardware. Figure~\ref{fig:cost} plots quality against both measures.
Against FLOPs, \fcsg{} sits far to the left of every other model. Against
latency it sits to the right of FFDNet with lower quality, so on a T4 FFDNet
beats it on both axes. We have not tested whether this changes on a CPU or a
mobile accelerator, where depthwise convolutions run better.

\begin{figure}[htbp]
	\centering
	\includegraphics[width=\linewidth]{compute_latency.pdf}
	\caption{Validation PSNR against reported GFLOPs (left) and against
	measured $256\times256$ forward latency (right).}
	\label{fig:cost}
\end{figure}

\paragraph{Tiling.} Tiled inference changes what each model sees. The FFT,
the gate's pooling windows and the global pooling in the fusion block all
work on a $256\times256$ tile instead of the whole image. To measure the
effect we took ten fixed validation images (0801, 0812, \dots, 0900) and
evaluated them with tiles of 128, 256 and 512 pixels and a fixed 32-pixel
overlap (Table~\ref{tab:tiling}).

\begin{table}[htbp]
	\centering
	\caption{Mean PSNR (dB) on ten validation images at three tile sizes, with 32-pixel overlap.}
	\label{tab:tiling}
	\begin{tabular}{@{}lrrr@{}}
		\toprule
		Method & Tile 128 & Tile 256 & Tile 512 \\
		\midrule
		\fcsg{}         & 26.622 & 26.659 & 26.672 \\
		DnCNN           & 27.330 & 27.341 & 27.343 \\
		FFDNet, blind   & 27.696 & 27.714 & 27.721 \\
		FFDNet, oracle  & 27.730 & 27.748 & 27.756 \\
		\bottomrule
	\end{tabular}
\end{table}

As a model with global operations, \fcsg{} reacts most to tile size, but not
by much. Moving from 256 to 512 pixel tiles gains 0.013\dB{}, against
0.002\dB{} for DnCNN and about 0.008\dB{} for FFDNet. The 0.6 to 1.0\dB{} gap
in Table~\ref{tab:main} is a hundred times larger, so tiling does not explain
it. The ten-image means are lower than the 100-image means because this
subset happens to hold harder images. Compare them only with each other.
Figure~\ref{fig:tiling} shows the paired differences.

\begin{figure}[htbp]
	\centering
	\includegraphics[width=\linewidth]{tiling.pdf}
	\caption{Tile-size sensitivity. PSNR change relative to 256-pixel tiles,
	and the geometry of tiles and overlaps.}
	\label{fig:tiling}
\end{figure}
% ===== end chapters/ch5_results.tex =====

\thispagestyle{empty}\cleardoublepage
% ===== begin chapters/ch6_routing.tex =====
\chapter{Routing analysis}\label{ch:routing}

Chapter~\ref{ch:results} measured how well \fcsg{} restores images. Here we
look at what its gate does. Every measurement uses the final 200,000-update
checkpoint from Chapter~\ref{ch:results}. We took ten fixed validation images
(0801, 0812, \dots, 0900) and degraded each with eight random seeds, for 80
samples in all. From each sample we cut the $256\times256$ centre crop after
degrading the whole image. The ten source images are the independent units.
The eight seeds per image change only the degradation.

\section{How the gate trains}\label{sec:routing-training}

The entropy term should keep all three experts busy early in training, then
fade so the gate can commit. Figure~\ref{fig:entropy} tracks the dense
routing entropy during training, summed over the three bands. It peaks at
$3\ln 3 = 3.296$ nats when every expert gets equal probability everywhere.

\begin{figure}[htbp]
	\centering
	\includegraphics[width=0.85\linewidth]{training_entropy.pdf}
	\caption{Dense routing entropy summed over bands during training. The
	maximum is $3\ln 3 \approx 3.296$ nats.}
	\label{fig:entropy}
\end{figure}

When training ends, the dense probabilities are still close to uniform. On
the diagnostic samples the mean dense entropy per band is 1.089, 1.096 and
1.097 nats for the low, mid and high bands, with a ceiling of
$\ln 3 = 1.099$. The sparse top-2 weights are just as close to an even split,
at 0.684, 0.691 and 0.692 nats against $\ln 2 = 0.693$. On their own, these
numbers suggest the gate learned almost nothing. The rest of the chapter
shows that reading is wrong. Small gaps between logits still decide
\emph{which} two experts get picked, and that choice matters.

\section{Spatial routing statistics}\label{sec:routing-stats}

Figure~\ref{fig:routing-summary} sums up the routing over the 80 samples.
The gate picks two experts at every pixel, so the selection fractions in
each band add up to two.

\begin{figure}[htbp]
	\centering
	\includegraphics[width=\linewidth]{routing_summary.pdf}
	\caption{Mean sparse weight and selection fraction per expert and band
	over 80 samples.}
	\label{fig:routing-summary}
\end{figure}

The high band shows the sharpest pattern. The gate picks the $3\times3$
expert at 99.98\% of its pixels. The larger experts share the second slot and
still carry a fair amount of weight on average, about 0.26 for the
$7\times7$ expert and 0.22 for the $5\times5$. So the small kernel does not
own the high band outright. It is always there, and the gate mixes in one of
the larger kernels depending on location. This fits the idea behind the
design, that fine detail wants a small support. The routing data alone,
though, cannot show that the $3\times3$ expert does anything specific to high
frequencies.

Figure~\ref{fig:routing-map} gives the routing maps for one image. The maps
follow the image's structure. Selection changes along object boundaries and
between smooth and textured areas. A top-2 choice flips the moment two logits
swap rank, so smooth probabilities underneath can still produce hard edges
between selected pairs.

\begin{figure}[htbp]
	\centering
	\includegraphics[width=\linewidth]{routing_0834.pdf}
	\caption{Expert weights for validation image 0834 (seed 33). Rows show the
	low, mid and high band. Columns show the degraded crop and the weights of
	the $7\times7$, $5\times5$ and $3\times3$ experts on a shared 0 to 1 scale.}
	\label{fig:routing-map}
\end{figure}

These maps show what the gate outputs. They do not show what each expert
contributes to a given restored structure. That would need attribution
experiments, and we did not run any.

\section{Routing and the degradation parameters}\label{sec:routing-corr}

Every sample records its degradation parameters, so we can ask whether the
routing reacts to them. For each image we subtract the eight-sample mean from
both the mean routing weights and the degradation parameters. The comparison
then happens within images, and differences in content between images do not
drive it. We compute Pearson correlations between the nine mean expert
weights and the four parameters.

Several correlations are strong. In the high band, the mean weight of the
$3\times3$ expert correlates at $r = +0.81$ with the noise level. The noisier
the input, the more the gate relies on the small kernel for high frequencies.
In the low band, the weight of the $5\times5$ expert correlates at $+0.80$
with noise and that of the $7\times7$ expert at $-0.73$. More noise moves
low-band weight from the largest kernel to the middle one. The mid band moves
the same way, with the $3\times3$ weight at $+0.77$. All six of the strongest
correlations involve noise. The strongest one with blur is $+0.42$ and the
strongest with JPEG quality is $-0.39$. In this sample, the gate reacts
mostly to noise. Figure~\ref{fig:corr} shows the full correlation matrix next
to the intervention results of the next section.

\begin{figure}[htbp]
	\centering
	\includegraphics[width=\linewidth]{correlations_interventions.pdf}
	\caption{Left, within-image Pearson correlations between mean routing
	weights and degradation parameters. Right, per-image PSNR change under the
	inference interventions of Table~\ref{tab:interventions}.}
	\label{fig:corr}
\end{figure}

We treat these results as exploratory. We computed 36 correlations on ten
source images, with no correction for multiple comparisons and no
image-level confidence intervals. All four parameters also change together
in every sample, so a correlation with noise cannot be cleanly pulled apart
from the other corruptions. A causal claim would need a controlled sweep that
moves one parameter at a time. What the data do show is that the gate's
output is not constant. It moves with the degradation, in the same direction,
across images.

\section{Inference-time interventions}\label{sec:interventions}

A correlation shows that the gate reacts to its input. It does not show that
the network needs that reaction. To test that, we froze the trained weights
and replaced the gate's output at inference in two ways.
\begin{itemize}
	\item \emph{Equal pair weights.} Keep the two experts the gate picks at
	each pixel, but give each a weight of one half.
	\item \emph{Uniform weights.} Ignore the gate and give all three experts a
	weight of one third everywhere.
\end{itemize}
Table~\ref{tab:interventions} gives the PSNR change on the ten images.

\begin{table}[htbp]
	\centering
	\caption{PSNR change when we override the trained gate at inference (10 images).}
	\label{tab:interventions}
	\begin{tabular}{@{}lrr@{}}
		\toprule
		Intervention & Mean $\Delta$PSNR & Min / max \\
		\midrule
		Equal weights on the selected pair    & $-0.040$\dB & $-0.137$ / $+0.028$ \\
		Uniform weights over all three experts & $-0.372$\dB & $-0.986$ / $+0.163$ \\
		\bottomrule
	\end{tabular}
\end{table}

The two interventions pull apart two halves of the gate's decision.
Flattening the weights inside the chosen pair costs only 0.04\dB{} on
average, so the exact mix hardly matters. Dropping the selection costs
0.37\dB{} on average and close to 1\dB{} on one image. The trained network
depends mainly on \emph{which} experts are combined at each pixel and much
less on how they are weighted. This also explains how the gate can matter
while its entropy sits near the maximum. A near-uniform distribution over
three experts still leaves one of them consistently left out.

The test has a limit. It measures how much \emph{this} trained network leans
on its gate. The experts and fusion stage trained alongside the gate and
adapted to its choices, so overriding the gate puts the model somewhere it
has never been. The result shows that the gate does real work inside the
trained model. It does not show that a model trained without a gate from the
start would do worse. That question needs a matched training ablation, which
we list as future work (Section~\ref{sec:future}).

\section{Discussion}\label{sec:discussion}

Read together, Chapters~\ref{ch:results} and~\ref{ch:routing} give a mixed
but consistent result. The routing works in the sense we set out to test. The
gate produces maps that vary across the image, its choices follow the noise
level in a steady direction, and the trained model loses up to a decibel when
we override them. The design was meant to produce exactly these behaviours.

The routing does not make the model competitive. With under a hundred
thousand parameters, \fcsg{} finishes 0.6 to 1.0\dB{} behind baselines six to
nine times its size, and on a GPU it is slower than all of them. We think
the quality gap comes down to capacity. The largest expert has fewer than
40,000 parameters, and the network spends much of its width processing three
bands that add back up to the original image. Blind FFDNet, the strongest
fair model, wins on a plain, dense architecture, not on extra information.

Our reading is this. Frequency-band routing is a cheap way to make a small
model's behaviour \emph{interpretable}. We can see where each kernel size is
used and tie that to the degradation. Within our budget, it does not on its
own make the model better or faster.
% ===== end chapters/ch6_routing.tex =====

\thispagestyle{empty}\cleardoublepage
% ===== begin chapters/ch7_conclusion.tex =====
\chapter{Conclusion and future work}\label{ch:conclusion}

\section{Summary of findings}\label{sec:summary}

We asked two questions. Can a very small restoration network split an image
into frequency bands and route each band, pixel by pixel, to convolutional
experts with different kernel sizes? And does the learned routing mean
anything? Here we go back over the four objectives of
Section~\ref{sec:objectives}.

\paragraph{O1, build and verify.} \fcsg{} has 97,088 parameters and costs
9.39 reported GFLOPs at $256\times256$, far inside the limits of five million
and ten. Before training it passed every check. Bands reconstruct exactly,
frequency units hold across image sizes, gradients flow, mixed precision runs
cleanly, the model overfits a single image, and checkpoints restore exactly
(Section~\ref{sec:verification}).

\paragraph{O2, trustworthy baselines.} With the author's weights, our FFDNet
code matches the published CBSD68 result to within 0.01\dB{}. Our DnCNN
training loop reaches 32.65\dB{} on DIV2K for Gaussian denoising
(Section~\ref{sec:baseline-repro}).

\paragraph{O3, matched comparison.} Under one composite degradation and one
training protocol, \fcsg{} reaches 27.24\dB{} PSNR, 0.753 SSIM and 0.396
LPIPS on 100 DIV2K validation images. That puts it 0.59\dB{} below DnCNN and
0.95\dB{} below blind FFDNet, and FFDNet beats it on all 100 images. It gets
there with 82.6\% fewer parameters than DnCNN and 88.6\% fewer than FFDNet.
Its low FLOP count did not make it fast. On a Tesla T4 it is the slowest of
the four (Chapter~\ref{ch:results}).

\paragraph{O4, the routing.} The learned gate varies across the image. It
picks the $3\times3$ expert almost everywhere in the high band and shifts its
weights with the noise level, with within-image correlations up to
$r=+0.81$. Overriding its choices at inference costs 0.37\dB{} on average.
Flattening its weights inside the chosen pair costs only 0.04\dB{}. The
trained model cares more about which experts are picked than about how they
are mixed (Chapter~\ref{ch:routing}).

Put briefly, frequency-band routing gave us a small model whose behaviour we
can inspect and relate to the degradation. Within our budget it did not give
us a model that beats simple convolutional baselines on quality or speed.

\section{Limitations}\label{sec:limitations}

These results only go so far, for the following reasons.

We trained each model once without fixing the seed, so we have no estimate
of run-to-run variance. The gap holds across images, which makes a lucky
seed an unlikely explanation. Small differences, such as the one between the
two FFDNet variants, cannot be interpreted. All results come from the DIV2K
validation split, which we also used to monitor training. No test set was
left untouched.

The degradation is synthetic. Real archival photographs fade, scratch,
collect dust, pick up colour casts and carry scanner artefacts, and our
pipeline models none of these. We did not test on real archival images, and
no person rated the outputs. How any of these models handles our data tells
us nothing yet about real archives.

The routing analysis rests on ten source images. Its correlations are
exploratory. The four degradation parameters move together, and we applied
no correction for the number of comparisons. The inference interventions
measure how much the trained model leans on its gate. They do not measure
what a gate adds during training.

The cost figures come from one GPU, the Tesla T4, at batch size 1 in FP32.
We did not profile individual operators, so our account of why \fcsg{} is
slow is reasoning, not measurement. We did not measure latency on CPUs or
mobile hardware, or peak training memory at batch 16.

Last, \fcsg{} runs all nine expert passes whatever the gate decides, so its
sparse routing saves no computation.

\section{Future work}\label{sec:future}

The first experiment to run is a matched training ablation. Train \fcsg{} and
an identical model with fixed uniform weights, from the same initialisation
and seed, and compare them on the full validation set, ideally with several
seeds each. Only this experiment can show whether learning a router helps, as
opposed to whether a trained model leans on the router it has. Four related
ablations would each test one design choice from Chapter~\ref{ch:method}:
\begin{itemize}
	\item three identical $3\times3$ experts in place of three kernel sizes,
	\item a single per-image gate in place of the spatial gate,
	\item plain concatenation in place of SE fusion, and
	\item hard masks in place of soft ones.
\end{itemize}

Second, sparse weights could become sparse computation. The trained model
cares mostly about which two experts are picked. Skipping the third expert in
tiles where the gate never selects it could cut running cost by up to a
third. An operator-level profile should come first. If memory traffic and
launch overhead dominate latency, removing arithmetic may not help much. Timing
the model on a CPU and on a mobile accelerator would show whether its low
FLOP count pays off on hardware that handles depthwise convolutions well.

Third, the capacity gap could be attacked head on by widening the experts.
The parameter budget is about fifty times what the model uses now. The
compute budget is nearly spent, though, so wider experts would need cheaper
expert evaluation to go with them, for example one unshuffled feature map
shared across bands. Giving the gate more context than a $16\times16$ window
may also help.

Finally, any claim about archival restoration needs three more things. These
are real archival photographs, a controlled sweep that separates the effect
of each corruption on the routing, and a held-out test set that was never
touched during development.
% ===== end chapters/ch7_conclusion.tex =====

\thispagestyle{empty}\cleardoublepage

% ---------------------------------------------------------------- back matter
\rhead{\textit{References}}
\lhead{}
\renewcommand\bibname{References}
\bibliographystyle{Ref/IEEEtran}   % CS/EE norm; the template lists it as an option
{
	\fontsize{10}{12}
	\selectfont
	\addcontentsline{toc}{chapter}{References}
	\bibliography{Ref/references}
}
\cleardoublepage

\appendix
% ===== begin appendix/appA_reproducibility.tex =====
\chapter{Reproducibility record}\label{app:repro}

\section{Code and environment}

The project repository at \url{https://github.com/whynotramaa/fcsg-capstone}
holds all code, configuration files, logs and evaluation outputs. Model code
is in \texttt{models/} and \texttt{src/fcsg\_net/}, training in
\texttt{training/train.py}, and evaluation and plotting in
\texttt{evaluation/}. A configuration file in \texttt{configs/} selects each
model. The four composite runs use \texttt{fcsg.toml},
\texttt{dncnn\_composite.toml}, \texttt{ffdnet\_blind.toml} and
\texttt{ffdnet\_composite.toml}. They differ only in the model and the noise
map.

We ran training, verification and diagnostics on Kaggle with an NVIDIA Tesla
T4, Python 3.12.13, PyTorch 2.10.0 with CUDA 12.8, NumPy 2.0.2 and Pillow
11.3.0. Local machines were only for editing and packaging code.

\section{Where each result comes from}

\begin{table}[htbp]
	\centering
	\caption{Source files for the results in this report.}
	\label{tab:sources}
	\small
	\begin{tabular}{@{}ll@{}}
		\toprule
		Result & File \\
		\midrule
		Verification checks (Table~\ref{tab:gates}) & \texttt{results/milestone\_run\_20260925/} \\
		Gaussian DnCNN (Table~\ref{tab:gaussian})   & \texttt{results/benchmark.csv} \\
		Composite results (Table~\ref{tab:main})     & \texttt{results/phase3/*/evaluation.json} \\
		Per-image results (Appendix~\ref{app:per-image}) & \texttt{results/phase3/per\_image\_scores.csv} \\
		Training curves                              & \texttt{results/phase3/*/train\_log.csv} \\
		Latency, tiling, routing                     & \texttt{results/diagnostics/} \\
		Figures                                      & \texttt{output/pdf/fcsg-report-assets/} \\
		\bottomrule
	\end{tabular}
\end{table}

\section{Artefact identifiers}

Each Kaggle run packed its source code into a bundle and recorded the
bundle's SHA-256 hash. Even after later edits, the hash pins down the exact
code behind each result. Table~\ref{tab:hashes} lists the first sixteen hex
digits of each hash. The source manifests hold the full values.

\begin{table}[htbp]
	\centering
	\caption{SHA-256 identifiers (first 16 hex digits) of the evaluated artefacts.}
	\label{tab:hashes}
	\small
	\begin{tabular}{@{}ll@{}}
		\toprule
		Artefact & SHA-256 prefix \\
		\midrule
		Composite training source bundle & \texttt{85d6e16f72ead013} \\
		Diagnostics source bundle        & \texttt{d473bc10b0ad0c1e} \\
		5,000-pair dataset manifest      & \texttt{7ce7f310bdb4d2d5} \\
		\fcsg{} checkpoint               & \texttt{8db0bb0f3c251ddc} \\
		DnCNN checkpoint                 & \texttt{afb0abfa022ede04} \\
		FFDNet blind checkpoint          & \texttt{981560c18cbea134} \\
		FFDNet oracle checkpoint         & \texttt{ef9560e8ed1ebae0} \\
		\bottomrule
	\end{tabular}
\end{table}

\section{How to reproduce}

\begin{enumerate}
	\item Attach the DIV2K dataset to a Kaggle notebook with a GPU. Run
	\texttt{notebooks/checks.ipynb}, then the bundled milestone notebook
	\texttt{notebooks/phase1\_phase2.ipynb}. Both stop at the first failed
	check.
	\item Build and push one training notebook per configuration with
	\texttt{python notebooks/build\_phase3.py fcsg ffdnet\_blind
	ffdnet\_composite dncnn\_composite}. To continue a run that outlasts one
	session, rebuild it with \texttt{--resume}. It picks up from its last
	checkpoint.
	\item At 200,000 updates, each run evaluates the 100 validation images and
	writes \texttt{evaluation.json} and \texttt{status.json}.
	\item To measure latency, tiling and routing, build the diagnostic notebook
	with \texttt{python notebooks/build\_diagnostics.py}. Draw the figures with
	\texttt{python evaluation/plots.py}.
\end{enumerate}
% ===== end appendix/appA_reproducibility.tex =====

\cleardoublepage
% ===== begin appendix/appB_per_image.tex =====
\chapter{Per-image results}\label{app:per-image}

Table~\ref{tab:per-image} gives, in dB, the input PSNR and each model's
output PSNR for all 100 DIV2K validation images. Bold marks the best output
for each image. The values come from
\texttt{results/phase3/per\_image\_scores.csv}, rounded to two decimal
places, so their means can differ from Table~\ref{tab:main} in the last
digit. \texttt{thesis/tools/per\_image\_table.py} generates the table.

{\small
\begin{longtable}{@{}lrrrrr@{}}
	\caption{Per-image PSNR (dB) on DIV2K validation.}\label{tab:per-image}\\
	\toprule
	Image & Input & DnCNN & FFDNet blind & FFDNet oracle & \fcsg{} \\
	\midrule
	\endfirsthead
	\toprule
	Image & Input & DnCNN & FFDNet blind & FFDNet oracle & \fcsg{} \\
	\midrule
	\endhead
	\bottomrule
	\endfoot
% ===== begin appendix/per_image_rows.tex =====
0801 & 28.09 & 28.22 & 29.37 & \textbf{29.39} & 29.32 \\
0802 & 27.37 & 33.95 & 33.93 & \textbf{34.12} & 32.99 \\
0803 & 15.71 & 33.77 & \textbf{34.48} & 34.47 & 32.08 \\
0804 & 16.16 & 26.95 & \textbf{27.06} & 27.03 & 26.43 \\
0805 & 14.82 & 25.82 & \textbf{25.97} & 25.95 & 25.64 \\
0806 & 18.14 & 26.57 & \textbf{26.82} & \textbf{26.82} & 26.08 \\
0807 & 18.11 & 20.51 & 20.63 & \textbf{20.65} & 20.15 \\
0808 & 16.25 & 25.80 & 25.92 & \textbf{25.93} & 25.33 \\
0809 & 22.81 & \textbf{31.57} & 31.52 & 31.56 & 31.48 \\
0810 & 16.54 & 25.48 & 25.61 & \textbf{25.63} & 25.26 \\
0811 & 15.96 & 26.43 & 26.58 & \textbf{26.61} & 26.17 \\
0812 & 17.27 & 28.60 & 28.87 & \textbf{28.91} & 27.94 \\
0813 & 28.36 & 31.28 & \textbf{31.86} & 31.84 & 31.12 \\
0814 & 15.68 & 26.52 & 26.73 & \textbf{26.74} & 26.02 \\
0815 & 17.65 & 29.74 & 30.01 & \textbf{30.05} & 29.16 \\
0816 & 25.01 & 30.11 & 30.49 & \textbf{30.54} & 29.70 \\
0817 & 29.35 & 31.27 & 31.97 & \textbf{32.10} & 30.91 \\
0818 & 17.80 & 26.89 & 27.14 & \textbf{27.16} & 26.17 \\
0819 & 25.06 & 28.17 & 28.18 & \textbf{28.20} & 27.89 \\
0820 & 20.31 & 26.03 & 26.22 & \textbf{26.23} & 24.74 \\
0821 & 27.80 & \textbf{31.93} & 31.79 & 31.89 & 29.79 \\
0822 & 17.92 & 27.18 & 27.42 & \textbf{27.47} & 26.53 \\
0823 & 26.21 & 27.79 & 28.18 & \textbf{28.26} & 27.23 \\
0824 & 25.07 & 27.88 & 28.48 & \textbf{28.51} & 26.91 \\
0825 & 21.15 & 27.46 & 27.59 & \textbf{27.64} & 26.56 \\
0826 & 25.15 & 27.29 & 27.15 & \textbf{27.52} & 25.76 \\
0827 & 30.95 & 34.00 & 34.11 & \textbf{34.15} & 32.97 \\
0828 & 17.10 & 18.76 & 18.61 & \textbf{18.83} & 17.93 \\
0829 & 15.75 & 25.43 & 25.50 & \textbf{25.52} & 25.16 \\
0830 & 17.82 & 26.30 & 26.44 & \textbf{26.52} & 25.62 \\
0831 & 29.44 & 29.73 & 33.24 & \textbf{33.46} & 31.39 \\
0832 & 20.16 & 27.15 & 27.54 & \textbf{27.61} & 26.22 \\
0833 & 25.07 & 33.64 & 34.12 & \textbf{34.17} & 32.75 \\
0834 & 16.05 & 24.65 & 24.86 & \textbf{24.90} & 24.38 \\
0835 & 23.60 & 24.43 & 24.51 & \textbf{24.55} & 24.20 \\
0836 & 15.53 & 25.50 & 25.79 & \textbf{25.81} & 24.62 \\
0837 & 15.65 & 25.92 & \textbf{26.10} & 26.08 & 25.12 \\
0838 & 29.15 & 36.36 & 36.85 & \textbf{36.92} & 36.42 \\
0839 & 17.69 & 28.06 & 28.19 & \textbf{28.23} & 27.62 \\
0840 & 27.16 & 29.08 & 29.07 & \textbf{29.11} & 28.72 \\
0841 & 17.67 & 25.90 & 26.09 & \textbf{26.15} & 25.44 \\
0842 & 26.95 & 26.16 & 28.87 & \textbf{28.91} & 28.23 \\
0843 & 20.44 & 35.45 & 36.51 & \textbf{37.02} & 33.54 \\
0844 & 36.25 & 41.00 & 41.39 & \textbf{41.43} & 38.57 \\
0845 & 23.08 & 26.39 & 26.67 & \textbf{26.75} & 24.37 \\
0846 & 16.09 & 24.52 & 24.95 & \textbf{25.00} & 23.35 \\
0847 & 23.40 & 25.60 & 26.37 & \textbf{26.44} & 25.23 \\
0848 & 21.99 & 25.80 & 26.04 & \textbf{26.10} & 25.11 \\
0849 & 19.78 & 22.61 & 22.72 & \textbf{22.82} & 21.66 \\
0850 & 21.91 & 28.95 & 29.10 & \textbf{29.15} & 28.08 \\
0851 & 20.99 & 26.99 & 27.24 & \textbf{27.26} & 26.38 \\
0852 & 26.22 & 28.52 & 28.64 & \textbf{28.69} & 28.10 \\
0853 & 19.11 & 31.31 & 31.76 & \textbf{31.86} & 30.17 \\
0854 & 22.68 & 26.73 & 26.85 & \textbf{26.90} & 25.98 \\
0855 & 28.26 & 30.76 & 30.97 & \textbf{31.14} & 30.76 \\
0856 & 21.16 & 23.49 & 23.72 & \textbf{23.76} & 23.11 \\
0857 & 13.93 & 32.51 & 32.84 & \textbf{32.89} & 31.57 \\
0858 & 16.78 & 26.08 & \textbf{26.22} & 26.20 & 25.84 \\
0859 & 23.45 & 26.47 & 26.57 & \textbf{26.59} & 26.21 \\
0860 & 20.09 & 21.46 & \textbf{21.79} & \textbf{21.79} & 21.16 \\
0861 & 23.50 & 26.18 & 26.43 & \textbf{26.61} & 24.14 \\
0862 & 16.49 & 31.67 & 31.84 & \textbf{31.85} & 31.37 \\
0863 & 17.92 & 30.76 & 30.97 & \textbf{31.04} & 30.30 \\
0864 & 26.45 & 29.89 & 30.00 & \textbf{30.03} & 29.58 \\
0865 & 25.90 & 26.64 & 26.88 & \textbf{26.94} & 26.44 \\
0866 & 20.34 & \textbf{26.77} & 26.67 & 26.76 & 26.43 \\
0867 & 25.37 & 27.57 & 28.14 & \textbf{28.16} & 26.71 \\
0868 & 19.78 & 32.57 & 32.95 & \textbf{33.01} & 31.52 \\
0869 & 16.42 & 23.58 & 23.70 & \textbf{23.71} & 23.38 \\
0870 & 15.03 & 26.23 & 26.35 & \textbf{26.38} & 25.92 \\
0871 & 26.18 & 30.48 & \textbf{31.07} & 31.01 & 30.26 \\
0872 & 16.17 & 24.13 & 24.22 & \textbf{24.25} & 23.71 \\
0873 & 22.14 & 23.43 & 23.70 & \textbf{23.79} & 23.12 \\
0874 & 22.38 & 30.14 & \textbf{30.21} & 30.17 & 29.50 \\
0875 & 15.92 & 22.67 & 22.81 & \textbf{22.84} & 22.12 \\
0876 & 18.52 & 22.90 & 22.99 & \textbf{23.06} & 22.49 \\
0877 & 17.99 & 33.85 & \textbf{34.22} & 34.20 & 31.10 \\
0878 & 22.03 & 30.80 & 31.10 & \textbf{31.14} & 29.44 \\
0879 & 21.48 & 27.56 & 27.73 & \textbf{27.78} & 26.76 \\
0880 & 14.46 & 28.77 & \textbf{29.01} & \textbf{29.01} & 28.40 \\
0881 & 25.37 & 25.34 & 27.01 & \textbf{27.07} & 26.27 \\
0882 & 15.19 & 30.54 & \textbf{30.90} & 30.89 & 29.44 \\
0883 & 21.42 & 25.12 & \textbf{25.35} & 25.32 & 24.63 \\
0884 & 18.18 & 26.04 & 26.20 & \textbf{26.25} & 25.42 \\
0885 & 15.06 & 21.95 & 22.08 & \textbf{22.12} & 21.67 \\
0886 & 19.79 & 34.28 & 34.55 & \textbf{34.63} & 33.56 \\
0887 & 20.01 & 24.85 & 24.99 & \textbf{25.02} & 24.35 \\
0888 & 25.69 & 29.76 & 30.10 & \textbf{30.13} & 29.19 \\
0889 & 17.31 & 28.44 & \textbf{28.57} & 28.54 & 28.06 \\
0890 & 17.79 & 24.69 & \textbf{24.78} & \textbf{24.78} & 24.46 \\
0891 & 20.27 & 26.51 & 26.73 & \textbf{26.77} & 25.32 \\
0892 & 15.73 & 27.55 & \textbf{27.83} & 27.82 & 26.36 \\
0893 & 24.56 & 31.73 & 31.87 & \textbf{31.94} & 31.10 \\
0894 & 22.61 & 28.85 & 29.21 & \textbf{29.25} & 28.07 \\
0895 & 15.54 & 21.17 & 21.21 & \textbf{21.25} & 20.98 \\
0896 & 28.23 & 35.13 & 36.53 & \textbf{36.60} & 34.86 \\
0897 & 17.12 & 21.44 & \textbf{21.52} & 21.49 & 21.11 \\
0898 & 19.19 & 32.31 & 32.57 & \textbf{32.60} & 31.20 \\
0899 & 25.23 & 24.41 & \textbf{28.13} & 28.09 & 26.20 \\
0900 & 19.34 & 27.46 & 27.65 & \textbf{27.67} & 26.03 \\

% ===== end appendix/per_image_rows.tex =====

\end{longtable}
}

% ===== end appendix/appB_per_image.tex =====

\cleardoublepage

% Dissemination: keep only if a paper/poster/preprint exists. Otherwise delete.
%\rhead{\textit{Dissemination}}
%\lhead{}
%\input{Ref/Dissemination}
%\thispagestyle{empty}
%\cleardoublepage

% Index: the template ends with \indexColumns{3}\indexPage. It is only useful
% if someone tags terms with \index{...} while writing. Drop it otherwise.
%\indexColumns{3}
%\indexPage
\end{document}
